% GENERATED FILE. Edit canonical lessons, then run npm run generate:books.
% canonical-source-hash: fnv1a64:b45d43f9300e4a7b
% canonical-lessons: JA-C155-hikkosu, JA-C155-hakobu, JA-C155-ugokasu, JA-C155-kagu, JA-C155-genkan

\chapter{Moving House and Shifting Furniture}
\label{ch:ja-moving-house-and-shifting-furniture}
\hlchaptermodality{155}

\begin{chapteropening}
\textbf{By the end of this chapter:} \emph{I can talk about moving into a home.}

Complete the last lesson of chapter 155: I can talk about moving into a home.
\end{chapteropening}

\section[\texorpdfstring{\ja{ひっこす} (hikkosu)}{ひっこす (hikkosu)}]{\ja{ひっこす} (hikkosu) — to move house}

\label{lesson:JA-C155-hikkosu}



\begin{quote}
Before the new one: say the Japanese for to get by, then the Japanese for to get used to.
\end{quote}

\subsection*{You'll want to know: \ja{ひっこす}}
\textbf{\ja{ひっこす}} — \emph{hikkosu} — \textquotedblleft{}to move house\textquotedblright{}.

\begin{quote}
\textbf{\ja{おじさんは} \ja{よく} \ja{ひっこします。}} — \emph{ojisan wa yoku hikkoshimasu} — \textquotedblleft{}My uncle often moves house.\textquotedblright{}
\end{quote}

Polite \textbf{\ja{ひっこします}}, \emph{hikkoshimasu}; a godan verb.

\begin{cousinweb}
\textbf{\ja{ひっこす}} (\emph{hikkosu}) is native Japanese, built from \emph{hiki}, \textquotedblleft{}pulling\textquotedblright{}, the stem of \emph{hiku}, and \emph{kosu}, \textquotedblleft{}to cross over\textquotedblright{}, joined with a held \emph{k}; both parts are inherited rather than borrowed (JMdict, KANJIDIC2).
\end{cousinweb}

\begin{grammarlens}[title={What you've built}]
One more word for everyday life at A2.
\end{grammarlens}

\subsection*{Your turn}
\begin{itemize}
\raggedright
  \item \emph{Say it:} \emph{hikkosu}
  \item \emph{Say it:} \emph{hikkosu}, once more
  \item \emph{Say it:} \emph{ojisan wa yoku hikkoshimasu} — My uncle often moves house.
  \item \emph{Recall:} say the Japanese for to get by, then the Japanese for to get used to, then say \emph{hikkosu} again
\end{itemize}

\begin{tcolorbox}[breakable,colback=teal!4,colframe=teal!35!black,title={Before you move on}]
What does \ja{ひっこす} mean? (\textquotedblleft{}To move house\textquotedblright{}.) Say it once more.
\end{tcolorbox}
\section[\texorpdfstring{\ja{はこぶ} (hakobu)}{はこぶ (hakobu)}]{\ja{はこぶ} (hakobu) — to carry}

\label{lesson:JA-C155-hakobu}



\begin{quote}
Before the new one: say the Japanese for to get used to, then the Japanese for to move house.
\end{quote}

\subsection*{You'll want to know: \ja{はこぶ}}
\textbf{\ja{はこぶ}} — \emph{hakobu} — \textquotedblleft{}to carry\textquotedblright{}.

\begin{quote}
\textbf{\ja{ちちと} \ja{にもつを} \ja{はこびます。}} — \emph{chichi to nimotsu o hakobimasu} — \textquotedblleft{}I carry the luggage with my father.\textquotedblright{}
\end{quote}

Polite \textbf{\ja{はこびます}}, \emph{hakobimasu}; a godan verb.

\begin{cousinweb}
\textbf{\ja{はこぶ}} (\emph{hakobu}) is a native Japanese word (wago), inherited rather than borrowed from Chinese; its kanji was attached later as a kun'yomi, and its older origin is not settled (JMdict, KANJIDIC2).
\end{cousinweb}

\begin{grammarlens}[title={What you've built}]
One more word for everyday life at A2.
\end{grammarlens}

\subsection*{Your turn}
\begin{itemize}
\raggedright
  \item \emph{Say it:} \emph{hakobu}
  \item \emph{Say it:} \emph{hakobu}, once more
  \item \emph{Say it:} \emph{chichi to nimotsu o hakobimasu} — I carry the luggage with my father.
  \item \emph{Recall:} say the Japanese for to get used to, then the Japanese for to move house, then say \emph{hakobu} again
\end{itemize}

\begin{tcolorbox}[breakable,colback=teal!4,colframe=teal!35!black,title={Before you move on}]
What does \ja{はこぶ} mean? (\textquotedblleft{}To carry\textquotedblright{}.) Say it once more.
\end{tcolorbox}
\section[\texorpdfstring{\ja{うごかす} (ugokasu)}{うごかす (ugokasu)}]{\ja{うごかす} (ugokasu) — to shift; to move (something)}

\label{lesson:JA-C155-ugokasu}



\begin{quote}
Before the new one: say the Japanese for to move house, then the Japanese for to carry.
\end{quote}

\subsection*{You'll want to know: \ja{うごかす}}
\textbf{\ja{うごかす}} — \emph{ugokasu} — \textquotedblleft{}to shift; to move (something)\textquotedblright{}.

\begin{quote}
\textbf{\ja{つくえを} \ja{うごかします。}} — \emph{tsukue o ugokashimasu} — \textquotedblleft{}I shift the desk.\textquotedblright{}
\end{quote}

Polite \textbf{\ja{うごかします}}, \emph{ugokashimasu}; a godan verb. Its intransitive partner is \textbf{\ja{うごく}}, \emph{ugoku}.

\begin{cousinweb}
\textbf{\ja{うごかす}} (\emph{ugokasu}) is a native Japanese word (wago), inherited rather than borrowed from Chinese, and built on the same root as \emph{ugoku}, \textquotedblleft{}to move\textquotedblright{}; its kanji was attached later as a kun'yomi (JMdict, KANJIDIC2).
\end{cousinweb}

\begin{grammarlens}[title={What you've built}]
One more word for everyday life at A2.
\end{grammarlens}

\subsection*{Your turn}
\begin{itemize}
\raggedright
  \item \emph{Say it:} \emph{ugokasu}
  \item \emph{Say it:} \emph{ugokasu}, once more
  \item \emph{Say it:} \emph{tsukue o ugokashimasu} — I shift the desk.
  \item \emph{Recall:} say the Japanese for to move house, then the Japanese for to carry, then say \emph{ugokasu} again
\end{itemize}

\begin{tcolorbox}[breakable,colback=teal!4,colframe=teal!35!black,title={Before you move on}]
What does \ja{うごかす} mean? (\textquotedblleft{}To shift; to move (something)\textquotedblright{}.) Say it once more.
\end{tcolorbox}
\section[\texorpdfstring{\ja{かぐ} (kagu)}{かぐ (kagu)}]{\ja{かぐ} (kagu) — furniture}

\label{lesson:JA-C155-kagu}



\begin{quote}
Before the new one: say the Japanese for to carry, then the Japanese for to shift.
\end{quote}

\subsection*{You'll want to know: \ja{かぐ}}
\textbf{\ja{かぐ}} — \emph{kagu} — \textquotedblleft{}furniture\textquotedblright{}.

\begin{quote}
\textbf{\ja{この} \ja{へやは} \ja{かぐが} \ja{すくないです。}} — \emph{kono heya wa kagu ga sukunai desu} — \textquotedblleft{}This room has little furniture.\textquotedblright{}
\end{quote}

\begin{cousinweb}
\textbf{\ja{かぐ}} (\emph{kagu}) is Sino-Japanese (kango): two Chinese characters read with their on'yomi, readings borrowed from Chinese. They are \emph{ka}, \textquotedblleft{}house, home\textquotedblright{} (Middle Chinese \emph{kae}) and \emph{gu}, \textquotedblleft{}tool, utensil\textquotedblright{} (\emph{gjuH}) (KANJIDIC2; Middle Chinese in Baxter's transcription).
\end{cousinweb}

\begin{grammarlens}[title={What you've built}]
One more word for everyday life at A2.
\end{grammarlens}

\subsection*{Your turn}
\begin{itemize}
\raggedright
  \item \emph{Say it:} \emph{kagu}
  \item \emph{Say it:} \emph{kagu}, once more
  \item \emph{Say it:} \emph{kono heya wa kagu ga sukunai desu} — This room has little furniture.
  \item \emph{Recall:} say the Japanese for to carry, then the Japanese for to shift, then say \emph{kagu} again
\end{itemize}

\begin{tcolorbox}[breakable,colback=teal!4,colframe=teal!35!black,title={Before you move on}]
What does \ja{かぐ} mean? (\textquotedblleft{}Furniture\textquotedblright{}.) Say it once more.
\end{tcolorbox}
\section[\texorpdfstring{\ja{げんかん} (genkan)}{げんかん (genkan)}]{\ja{げんかん} (genkan) — the front entrance; the entry hall}

\label{lesson:JA-C155-genkan}



\begin{quote}
Before the new one: say the Japanese for to shift, then the Japanese for furniture.
\end{quote}

\subsection*{You'll want to know: \ja{げんかん}}
\textbf{\ja{げんかん}} — \emph{genkan} — \textquotedblleft{}the front entrance; the entry hall\textquotedblright{}.

\begin{quote}
\textbf{\ja{げんかんで} \ja{くつを} \ja{ぬぎます。}} — \emph{genkan de kutsu o nugimasu} — \textquotedblleft{}I take off my shoes in the entry hall.\textquotedblright{}
\end{quote}

\begin{cousinweb}
\textbf{\ja{げんかん}} (\emph{genkan}) is Sino-Japanese (kango): two Chinese characters read with their on'yomi, readings borrowed from Chinese. They are \emph{gen}, \textquotedblleft{}dark, deep\textquotedblright{} (Middle Chinese \emph{hwen}) and \emph{kan}, \textquotedblleft{}barrier, gateway\textquotedblright{} (\emph{kwaen}) (KANJIDIC2; Middle Chinese in Baxter's transcription).
\end{cousinweb}

\begin{grammarlens}[title={What you've built}]
One more word for everyday life at A2.
\end{grammarlens}

\subsection*{Your turn}
\begin{itemize}
\raggedright
  \item \emph{Say it:} \emph{genkan}
  \item \emph{Say it:} \emph{genkan}, once more
  \item \emph{Say it:} \emph{genkan de kutsu o nugimasu} — I take off my shoes in the entry hall.
  \item \emph{Recall:} say the Japanese for to shift, then the Japanese for furniture, then say \emph{genkan} again
\end{itemize}

\begin{tcolorbox}[breakable,colback=teal!4,colframe=teal!35!black,title={Before you move on}]
What does \ja{げんかん} mean? (\textquotedblleft{}The front entrance; the entry hall\textquotedblright{}.) Say it once more.
\end{tcolorbox}
